\documentclass{article}
\usepackage{amsmath, amssymb, amsthm}
\usepackage{hyperref}
\usepackage{geometry}
\geometry{margin=1in}

\title{Erd\H{o}s Problem \#685}
\date{Last updated: 2025-08-31}
\author{Source: erdosproblems.com}

\begin{document}
\maketitle

\noindent\textbf{Status:} OPEN \quad \textbf{No prize}

\noindent\textbf{Tags:} number theory, primes, binomial coefficients

\medskip

\noindent\textbf{Problem Statement:}

Let $\epsilon&#62;0$ and $n$ be large depending on $\epsilon$. Is it true that for all $n^\epsilon&#60;k\leq n^{1-\epsilon}$ the number of distinct prime divisors of $\binom{n}{k}$ is\[(1+o(1))k\sum_{k&#60;p&#60;n}\frac{1}{p}?\]Or perhaps even when $k \geq (\log n)^c$?




It is trivial that the number of prime factors is\[&#62;\frac{\log \binom{n}{k}}{\log n},\]and this inequality becomes (asymptotic) equality if $k&#62;n^{1-o(1)}$.


\bigskip
\noindent\small{Source: \url{https://www.erdosproblems.com/685}}
\end{document}
